%% kapa: exam=1 round=33 year=2022 period=1 form=A subject=real_estate subject_ko=부동산학원론 seq=33 number=113
%% source: Q-Net 감정평가사 기출문제 (cst00309) artlSeq=5209997 file=raw/1008/5209997/2022년도 제33회 감정평가사 1차시험 1교시(민법,경제학,부동산학).pdf
%% conversion: kordoc 4.24.51 -> tools/gen1.py (auto); review=proofread-sonnet; flags=none
\begin{question}{113}
부동산 권리분석에 관한 설명으로 옳지 않은 것은?
\begin{choices}
  \item 권리분석의 원칙에는 능률성, 안전성, 탐문주의, 증거주의 등이 있다.
  \item 건물의 소재지, 구조, 용도 등의 사실관계는 건축물대장으로 확인·판단한다.
  \item 임장활동 이전 단계 활동으로 여러 가지 물적 증거를 수집하고 탁상으로 검토하여 1차적으로 하자의 유무를 발견하는 작업을 권리보증이라고 한다.
  \item 부동산의 상태 또는 사실관계, 등기능력이 없는 권리 및 등기를 요하지 않는 권리관계 등 자세한 내용까지 분석의 대상으로 하는 것이 최광의의 권리분석이다.
  \item 매수인이 대상부동산을 매수하기 전에 소유권을 저해하는 조세체납, 계약상 하자 등을 확인하기 위해 공부 등을 조사하는 일도 포함된다.
\end{choices}
\end{question}
\answer{3}
